\chapter{Introduction}
\label{ch:intro}

LiNC is the operator interface of a LinuxCNC machine. It shows what the machine is
doing, it loads and edits G-code, and it holds the tool, offset and probing pages.
It runs on the control PC as a single program: there is no helper process, no Python
layer and no web service between the screen and the controller.

LiNC is operated from three surfaces, and none of them is a touch screen:

\begin{enumerate}
  \item The \textbf{operator panel} -- power, emergency stop, mode selector, override
        knobs, jog keys, cycle keys and the machine function keys. This is where the
        machine is controlled.
  \item The \textbf{CNC keyboard} -- navigation, G-code entry, MDI, tool and offset
        values, file handling. Every action on the screen has a key.
  \item The \textbf{screen} -- readouts, the toolpath, the program, the tables and the
        diagnostics. It reports; it does not control motion.
\end{enumerate}

The layout and the key handling follow the conventions of an industrial machine
controller, so an operator who knows a conventional control can work without learning
a new interaction model. A mouse is not required, and a touch screen is not required.

The panel and the interface can also run on one computer without a machine cabinet: the
\emph{simulated panel} is a second window that carries the same keys, knobs, lamps and
displays, and publishes the same interface to the controller. On a machine without a
physical panel it stands in for one; a touch screen or a mouse then works as a
replacement for the panel keys. It is a replacement for a panel, not a second way to
drive the machine.

Because the machine controls stay off the screen, the screen is free for machine state.
The digital readout, the toolpath, the program listing, the modal codes and the servo
load are visible at the same time, which is what supervising a running job needs.

\shot{img/ui/machine.png}{The Machine page with a program loaded: readout, toolpath,
program listing, feed and spindle, servo load.}{fig:intro-machine}

\section{What makes LiNC different}

\begin{description}[font=\normalfont\sffamily\bfseries, leftmargin=4.6cm, style=nextline]
  \item[The panel stays the panel] There is no on-screen jog, no on-screen cycle start,
        no override slider and no software emergency stop. A second copy of a control
        drifts out of agreement with the hardware, so there is no second copy.
  \item[MDI is a working area] The MDI pane keeps the blocks you have issued and draws
        the block you are typing on the toolpath before it is sent. A block that would
        leave the machine's travel is refused by name.
  \item[The toolpath explains itself] The loaded program, the live tool trail, the
        machine envelope and the part bounds are separate layers with separate colours.
        Anything outside the machine's travel is drawn solid red.
  \item[Tool changes without MDI] Select a row in the tool table and press \keycap{T}.
        The change and its length compensation are issued for you.
  \item[The editor checks the work] The Program page is a G-code editor with syntax
        highlighting, undo, redo and a live preview. An edit reaches the drawing only
        after it parses.
  \item[Probing is guided] Categories, illustrated operations, parameters and results,
        using the Probe Basic G-code routines that ship with LiNC.
  \item[Help is on the machine] A searchable G-code reference, offline, readable while a
        program runs.
  \item[The network can be taken away] When a program is loaded or MDI is selected, LiNC
        can hold the machine's external network link down, so WiFi latency cannot reach
        the realtime path. \chordtwo{Ctrl}{U} unloads the program and releases it.
\end{description}

\section{What the software will not do}

These are design rules, not preferences. Each one exists so that the screen and the
machine cannot disagree.

\begin{itemize}
  \item It will not jog, home, start the spindle, change coolant or start a cycle.
        Those are panel actions, wired through the machine's HAL.
  \item It will not write an arbitrary HAL pin. It writes seven declared outputs, and
        the list is fixed in the profile and in the program.
  \item It will not save G-code, write a table, import a file or apply a probe result
        while the panel's \pbtn{PROGRAM PROTECT} switch is on.
  \item It will not home an axis. It reports homing and explains the sequence.
  \item It will not guess. A stale or invalid panel state is shown as unknown and fails
        closed.
\end{itemize}

\section{How to read this manual}

\begin{description}[font=\normalfont\sffamily\bfseries, leftmargin=4.2cm, style=nextline]
  \item[Chapter \ref{ch:power}] Powering up and down, emergency stop, protect switch,
        the network interlock.
  \item[Chapter \ref{ch:panel}] Every control on the operator panel.
  \item[Chapter \ref{ch:keyboard}] The keyboard, the navigation column and every
        shortcut.
  \item[Chapter \ref{ch:screen}] The fixed shell: top bar, bottom bar, banners,
        messages.
  \item[Chapter \ref{ch:modes}] Panel modes, jogging, homing and program zero.
  \item[Chapter \ref{ch:running}] Loading, checking and running a program.
  \item[Chapter \ref{ch:mdi}] MDI and the G-code editor.
  \item[Chapter \ref{ch:tools}] Tool table, tool changes, offsets.
  \item[Chapter \ref{ch:probing}] Touch probing.
  \item[Chapter \ref{ch:diagnostics}] Diagnostics, logs and troubleshooting.
  \item[Chapter \ref{ch:quickref}] Key tables and reference data.
\end{description}

Boxes used throughout the manual:

\begin{safetybox}
  Instructions that prevent injury. Read them before operating the machine.
\end{safetybox}
\begin{cautionbox}
  Instructions that prevent damage to the machine, the tooling or the workpiece.
\end{cautionbox}
\begin{notebox}
  Context and cross references.
\end{notebox}
\begin{tipbox}
  Practical hints from the shop floor.
\end{tipbox}

\section{Notation}

Panel keys are drawn as \pbtn{CYCLE START}, screen actions and labels as \sbtn{SAVE},
keyboard keys as \keycap{F6} and key chords as \chordtwo{Ctrl}{S}. Lit panel
lamps are drawn as \ledlit{ledgreen} (lit) against \leddim{ledgreen} (dark). A state
word the screen paints is written as \state{NOT HOMED}.

Positions on the millimetre machine are in millimetres, speeds in millimetres per
minute, spindle speed in revolutions per minute. A lathe profile shows X in the
radius or diameter convention its INI selects.

\section{Terms used in this manual}

\begin{longtable}{L{3.4cm}L{9.6cm}}
  \toprule
  \textbf{Term} & \textbf{Meaning} \\
  \midrule
  \endhead
  HMI & Human-machine interface: this program. \\
  Panel & The operator panel, physical or simulated. It owns the machine controls. \\
  Controller & LinuxCNC: task, motion and I/O. It is authoritative for machine state. \\
  Task mode & LinuxCNC's own mode (MANUAL, MDI, AUTO). The panel requests it; LiNC
              displays it. \\
  Panel mode & The position of the mode selector on the panel. LiNC reads it as one-hot
              signals and fails closed if it is invalid. \\
  WCS & Work coordinate system, \texttt{G54} to \texttt{G59.3}. \\
  DTG & Distance to go: how far the tool still has to move. \\
  MDI & Manual data input: one block, issued immediately. \\
  HAL & LinuxCNC's hardware abstraction layer. \\
  Protect & The panel's \pbtn{PROGRAM PROTECT} key switch. It locks every write. \\
  Interlock & The rule that holds the external network link down while the machine can
              move. \\
  \bottomrule
\end{longtable}
